\documentclass[10pt,a4paper]{amsart}
\usepackage{amsmath,amssymb,amsthm,graphicx}
% For commutative diagrams we use Xy-pic
\usepackage[all]{xy}

% We use 2cell for 2-commutative diagrams.
\xyoption{2cell}
% Theorem environments.
%
\theoremstyle{plain}
\newtheorem{theorem}[subsection]{Theorem}
\newtheorem{proposition}[subsection]{Proposition}
\newtheorem{lemma}[subsection]{Lemma}

\theoremstyle{definition}
\newtheorem{definition}[subsection]{Definition}
\newtheorem{example}[subsection]{Example}
\newtheorem{exercise}[subsection]{Exercise}
\newtheorem{situation}[subsection]{Situation}

\theoremstyle{remark}
\newtheorem{remark}[subsection]{Remark}
\newtheorem{remarks}[subsection]{Remarks}

\numberwithin{equation}{subsection}

% Macros
%
\def\lim{\mathop{\mathrm{lim}}\nolimits}
\def\colim{\mathop{\mathrm{colim}}\nolimits}
\def\Spec{\mathop{\mathrm{Spec}}}
\def\Hom{\mathop{\mathrm{Hom}}\nolimits}
\def\Ext{\mathop{\mathrm{Ext}}\nolimits}
\def\SheafHom{\mathop{\mathcal{H}\!\mathit{om}}\nolimits}
\def\SheafExt{\mathop{\mathcal{E}\!\mathit{xt}}\nolimits}
\def\Sch{\mathit{Sch}}
\def\Mor{\mathop{\mathrm{Mor}}\nolimits}
\def\Ob{\mathop{\mathrm{Ob}}\nolimits}
\def\Sh{\mathop{\mathit{Sh}}\nolimits}
\def\NL{\mathop{N\!L}\nolimits}
\def\CH{\mathop{\mathrm{CH}}\nolimits}
\def\proetale{{pro\text{-}\acute{e}tale}}
\def\etale{{\acute{e}tale}}
\def\QCoh{\mathit{QCoh}}
\def\Ker{\mathop{\mathrm{Ker}}}
\def\Im{\mathop{\mathrm{Im}}}
\def\Coker{\mathop{\mathrm{Coker}}}
\def\Coim{\mathop{\mathrm{Coim}}}

% Boxtimes
%
\DeclareMathSymbol{\boxtimes}{\mathbin}{AMSa}{"02}

%
% Macros for moduli stacks/spaces
%
\def\QCohstack{\mathcal{QC}\!\mathit{oh}}
\def\Cohstack{\mathcal{C}\!\mathit{oh}}
\def\Spacesstack{\mathcal{S}\!\mathit{paces}}
\def\Quotfunctor{\mathrm{Quot}}
\def\Hilbfunctor{\mathrm{Hilb}}
\def\Curvesstack{\mathcal{C}\!\mathit{urves}}
\def\Polarizedstack{\mathcal{P}\!\mathit{olarized}}
\def\Complexesstack{\mathcal{C}\!\mathit{omplexes}}
% \Pic is the operator that assigns to X its picard group, usage \Pic(X)
% \Picardstack_{X/B} denotes the Picard stack of X over B
% \Picardfunctor_{X/B} denotes the Picard functor of X over B
\def\Pic{\mathop{\mathrm{Pic}}\nolimits}
\def\Picardstack{\mathcal{P}\!\mathit{ic}}
\def\Picardfunctor{\mathrm{Pic}}
\def\Deformationcategory{\mathcal{D}\!\mathit{ef}}
\usepackage{cleveref}
\crefname{originallemma}{original lemma}{original lemmas}
\Crefname{originallemma}{Original lemma}{Original lemmas}
\makeatletter
\newlabel{lemma-gysin-at-infty}{{0F9J}{source}{}{sourcepointer-1}{}}
\newlabel{lemma-gysin-at-infty@cref}{{[originallemma][0][]0F9J}{[1][0][]source}{}{}{}}
\makeatother
\crefname{originallemma}{original lemma}{original lemmas}
\Crefname{originallemma}{Original lemma}{Original lemmas}
\makeatletter
\newlabel{lemma-closed-in-X-gysin}{{02TA}{source}{}{sourcepointer-2}{}}
\newlabel{lemma-closed-in-X-gysin@cref}{{[originallemma][0][]02TA}{[1][0][]source}{}{}{}}
\makeatother
\title{Independence of the Gysin-at-infinity class from the choice of the proper model, up to the induced morphism at infinity}
\author{The Stacks Project Authors}
\date{stacks-project:chow.tex:local-snapshot-20261012; GNU FDL 1.2 or later}
\begin{document}
\maketitle
\noindent\textit{Editorial note.} All delivered mathematical text is the author's own, copied from the pinned Stacks Project chapter file chow.tex; no statement or proof step was rewritten, reordered or removed.\par\medskip
\noindent\textit{Editorial note.} Preamble provenance: the pinned chapter file opens with an input of the shared preamble and carries no preamble text of its own, so the project's single Stacks-wide preamble staged at sources/stacks-chow/preamble.tex is excerpted verbatim as two spans, lines 14-18 (Xy-pic setup) and lines 154-217 (the native theorem declarations and mathematical macros). Nothing was added, removed or redefined; the wrapper class and the amsmath/amssymb/amsthm/graphicx base packages are supplied by the assembler.\par\medskip
\noindent\textit{Editorial note.} The counted statement and proof contain no cross reference, no citation and no included file, so the delivered body is self contained.\par\medskip
\noindent\textit{Editorial note.} Canonical location of the counted proposition: Stacks tag 0FAQ (sources/stacks-chow/pinned-original-source.tex line 10108).\par\medskip
\noindent\textit{Allowed standard prerequisites (editorial).}\begin{itemize}
\item The Gysin-at-infinity bivariant class and its defining property. Same chapter, Lemma labelled lemma-gysin-at-infty (source line 10026), invoked in the statement. Supplies the two classes C and C' and the characterisation C cap alpha = i\_infinity\textasciicircum{}* beta used to compare them.
\item Compatibility of the Gysin homomorphism with proper pushforward. Same chapter, Lemma labelled lemma-closed-in-X-gysin (source line 5264), invoked at proof line 10125. Supplies i\_infinity\textasciicircum{}* beta = g\_\{infinity,*\}(i'\_infinity)\textasciicircum{}* beta', which is the identity that makes the two constructions agree.
\end{itemize}\medskip
\ifdefined\hypertarget\hypertarget{sourcepointer-1}{}\fi
\noindent\textit{Source reference.} Original lemma 0F9J: sources/stacks-chow/pinned-original-source.tex, source lines 10026--10026. Same-chapter pointer into the author environment lemma at source line 10026; invoked at proof line None of the counted proof. It is retained only as a navigation pointer; every load-bearing step is inside the delivered span.\par\medskip
\ifdefined\hypertarget\hypertarget{sourcepointer-2}{}\fi
\noindent\textit{Source reference.} Original lemma 02TA: sources/stacks-chow/pinned-original-source.tex, source lines 5264--5264. Same-chapter pointer into the author environment lemma at source line 5264; invoked at proof line 10125 of the counted proof. It is retained only as a navigation pointer; every load-bearing step is inside the delivered span.\par\medskip
\begin{lemma}
\label{lemma-gysin-at-infty-independent}
In Lemma \ref{lemma-gysin-at-infty} let $g : W' \to W$ be a proper morphism
which is an isomorphism over $\mathbf{A}^1_X$. Let
$C' \in A^0(W'_\infty \to X)$ and $C \in A^0(W_\infty \to X)$
be the classes constructed in Lemma \ref{lemma-gysin-at-infty}.
Then $g_{\infty, *} \circ C' = C$ in $A^0(W_\infty \to X)$.
\end{lemma}

\begin{proof}
Set $b' = b \circ g : W' \to \mathbf{P}^1_X$. Denote
$i'_\infty : W'_\infty \to W'$ the inclusion morphism.
Denote $g_\infty : W'_\infty \to W_\infty$ the restriction of $g$.
Given $\alpha \in \CH_k(X)$ choose $\beta' \in \CH_{k + 1}(W')$
restricting to the flat pullback of $\alpha$ on $(b')^{-1}\mathbf{A}^1_X$.
Then $\beta = g_*\beta' \in \CH_{k + 1}(W)$ restricts to the
flat pullback of $\alpha$ on $b^{-1}\mathbf{A}^1_X$.
Then $i_\infty^*\beta = g_{\infty, *}(i'_\infty)^*\beta'$
by Lemma \ref{lemma-closed-in-X-gysin}.
This and the corresponding fact after base change by
morphisms $X' \to X$ locally of finite type, corresponds
to the assertion made in the lemma.
\end{proof}

\end{document}
